% Place after the bibliography.
% Required packages: booktabs, array.
% Replace the previous appendix block; do not append this block to it.
% Use \appendix only once in each language's manuscript.

\clearpage
\appendix

% Appendix A
\section{Reproducibility and Measurement}
\label{app:protocol}

\paragraph{Evidence binding.} The machine-readable package in `appendix\_evidence/` binds the final 96-row export, cohort manifests, source hashes, and post-processing commands. Table sources are `TABLE\_LATENCY.csv`, `TABLE\_E2\_MEMORY.csv`, and `TABLE\_HISTORICAL\_PAIRS.csv`; code correspondence is in `STATE\_CODE\_MAP.csv`. Missing or unevaluable fields are recorded in `DATA\_GAPS.md`.


The cold-resumption study uses the same Lychee-FD model and local
Token2Wav backend \citep{lychee2026} on two A100-SXM4-40GB GPUs,
with model execution and one acoustic lane placed separately.
Shared weights remain resident.
Model, sampling, numerical, input, and playback settings are matched
within each comparison. Reproduction requires the executed source,
dependency, and configuration manifests, including archived patches,
rather than a repository commit alone.

The final matrix contains 96 executions from 12 source groups,
four policies, and two repetitions. E1 contains 44 executions and
E2 contains 52; their measurements are analyzed separately.
The 92 validated cuts contain 24 pending speech tokens and
20,032 buffered PCM samples at 24\,kHz.
These values describe the predefined nonempty cut, not a natural
conversation-state distribution.
The nominal suspension is 30 seconds, but latency uses actual events.
New-audio rendering spans Worklet resumption to the first-new-PCM frame
rendered within the same AudioContext;
buffered playback is not new computation.

Configuration summaries aggregate valid repetitions within each
source and environment before taking the median across sources.
Paired differences first match source, repetition, and environment.
Missing values are not replaced with zero.
The four unvalidated executions comprise three frame-timeline safety
cancellations and one cut-alignment rejection; they remain in the
reported denominators.
The cold, transfer, throughput, and interaction cohorts retain
separate identities and are not pooled.


% Appendix B
\section{Joint State and Resumption Validation}
\label{app:state_contract}

Table~\ref{tab:app_state_contract} summarizes the joint resumption
boundary. Computational state is restored, whereas the external ledger
of accepted input, cancellation decisions, and delivery records remains current.
This distinction applies consistent-snapshot and rollback principles at request granularity \citep{chandy1985,elnozahy2002}.

\begin{table}[tbp]
    \centering
    \small
    \setlength{\tabcolsep}{5pt}
    \caption{
        State coverage and resumption obligations.
        External records are retained rather than overwritten
        with an older checkpoint copy.
    }
    \label{tab:app_state_contract}
    \begin{tabular}{
        @{}
        >{\raggedright\arraybackslash}p{0.20\linewidth}
        >{\raggedright\arraybackslash}p{0.73\linewidth}
        @{}
    }
        \toprule
        State class & Contents and resumption obligation \\
        \midrule
        Model
        & Effective KV, non-KV histories, processor/control state,
          sampling context, and logical positions.
          Restore values in legal mappings and rebuild active bindings. \\
        Bridge
        & Produced but unconsumed tokens, partial buffers,
          initialization/flush state, and pending output.
          Preserve content and consumption order. \\
        Acoustic
        & StreamingDecoder and Token2Wav continuation,
          Flow / HiFT caches, actual random context, and pending PCM.
          Restore a compatible execution context. \\
        Identity and progress
        & Request, stream, generation, sequence, state version,
          and stage watermarks.
          Validate ownership against the current lifecycle. \\
        External records
        & Accepted input, cancellation, and delivery/rendering intervals.
          Preserve current facts and reject obsolete submissions. \\
        \bottomrule
    \end{tabular}
\end{table}

Capture establishes a consistent cut and an independent host
checkpoint before releasing request-private resources.
Resumption prepares legal mappings, rebuilds bindings, and validates
the participating states before enabling execution or output.
Current generation, version, and cancellation state are rechecked
at publication; a failed or cancelled resumption is not published.
The restored input-consumption position does not overwrite input
received after the checkpoint.

The request ID regression distinguishes physical output order
$[A,B]$ from a sampler selection containing only $B$.
request ID lookup retrieves $B$'s corresponding side outputs rather
than physical row zero.
Separate archived checks cover pending bridge content,
pre-publication rejection, and request-local cancellation.
Complete token and waveform comparisons are conditioned on matched
input/control histories and numerical configurations; they do not
establish semantic correctness or unconditional kernel determinism.


% Appendix C
\section{Additional Performance Results}
\label{app:performance}

Table~\ref{tab:app_latency} supplements the main rendering results.
Save measures pause acceptance to independent-checkpoint readiness; Ready
measures RESUME acceptance to execution readiness. Gap measures buffered PCM
exhaustion to the first-new-PCM frame rendered by the AudioContext.
Each interval retains its own endpoints and is not added to other
intervals to construct a new end-to-end metric.

\begin{table}[tbp]
    \centering
    \small
    \setlength{\tabcolsep}{6pt}
    \caption{
        Additional cold-resumption measurements in seconds.
        Values are source-level medians within each environment.
        Valid/all retains the complete scheduled denominator.
    }
    \label{tab:app_latency}
    \begin{tabular}{@{}llcrrr@{}}
        \toprule
        Environment & Policy & Valid/all & Save & Ready & Gap \\
        \midrule
        E1 & Hot     & 11/11 & 0.004 & 0.001  & 0.349 \\
        E1 & Rebuild & 10/11 & 1.847 & 13.317 & 13.616 \\
        E1 & Joint   & 11/11 & 1.387 & 0.699  & 1.064 \\
        E1 & Hybrid  & 8/11  & 0.655 & 0.130  & 0.515 \\
        \midrule
        E2 & Hot     & 13/13 & 0.004 & 0.001  & 0.333 \\
        E2 & Rebuild & 13/13 & 1.926 & 13.250 & 13.565 \\
        E2 & Joint   & 13/13 & 1.387 & 0.700  & 1.037 \\
        E2 & Hybrid  & 13/13 & 0.696 & 0.136  & 0.480 \\
        \bottomrule
    \end{tabular}
\end{table}

In the final cohort, rendering wait exceeds the post-buffer gap
by $20{,}032/24{,}000=0.834667$ seconds.
These are related views of the same delivery timeline, not
independent improvements.
In E2, Joint and Hybrid hold independent checkpoint payloads of
672.207 and 22.416\,MiB, respectively, and managed pinned buffers
of 256 and 128\,MiB.
Both retain model-device workspaces, producing a net allocation
increase of 60.489\,MiB.
Checkpoint size, pinned storage, process RSS, and device allocation
are distinct, non-additive measurements.

The independent transfer study contains 12 Sync/Optimized pairs
from six source groups.
Relative reduction is calculated within each pair as
$(T_i^{\mathrm{sync}}-T_i^{\mathrm{opt}})/T_i^{\mathrm{sync}}$.
Its median is 16.00\%, and the median absolute reduction is
0.533 seconds, with all 12 pairs favoring the optimized path.
The additional paused and sampled peak RSS are median paired
243.2 and 182.2\,MiB, respectively.
The archived throughput matrix instead averages within-repeat
ratios for each workload/concurrency cell.
Its full set of configurations and repetitions must be retained;
new throughput measurements form a separate cohort.


% Appendix D
\section{Interaction and External Protocols}
\label{app:interaction}

The natural-interaction pilot uses adapted FD-Bench-style definitions
\citep{fdbench2025}.
SRR retains eligible non-interruption questions, SIR actual
interruption opportunities, and SRIR successfully interrupted events
as their respective denominators.
EIR records premature assistant entry.
FSED spans user-utterance end to matched reply playback onset;
IRD spans interruption onset to the end of preceding assistant speech.
Signed timings are retained.
Explicit CPU suspension is not counted as a speech interruption,
and PCM concatenation does not replace the actual rendering timeline.

The shared Lychee/DuplexPilot online path is one implementation,
not two independent baselines.
It uses 400\,ms end-of-chunk input delivery, whereas the Omni client
uses 200\,ms send-then-wait delivery.
Models and actual interruption exposure also differ.
The pilot therefore describes application behavior rather than
isolating planned-cold-resumption effects.
Prior listening feedback applies only to the reviewed examples and
is not extended to unreviewed pilot outputs.

Two native vLLM-Omni hot reattachments \citep{omnidoc} produce
request-to-ACK times of 1.512 and 1.421\,ms and first-PCM reception
times of 725.412 and 776.220\,ms.
Recorded PCM delivery has no missing or duplicate sequence events
in these cases, but the production phase of the first returned PCM
is unknown.
The LiveServe-style preloading component \citep{liveserve2026}
prepares 20 and 15\,MiB of KV approximately 9.9 seconds in advance.
It improves readiness by approximately 56.65 and 100.01\,ms,
while new-PCM reception from the common source-window anchor is
later by 23.7 and 36.4\,ms.
These observations retain their native protocols and do not supply
a unified external planned-cold-resumption comparison.